\begin{lemma}[The principal value of a truncated Gaussian Riesz integral exists almost everywhere]
Fix $n \geq 1$, let $\lambda$ be Lebesgue measure on $\mathbb{R}^n$, let
$\gamma_n$ be the Gaussian measure of \noderef{GaussMeasure}, and let
$K_{\gamma,j}$ be the kernel of \noderef{GaussRieszKernel}.  Let
\[
  T : L^1(\gamma_n;\mathbb{R}) \longrightarrow
      \mathcal{M}\bigl(\gamma_n;\mathbb{R}^n\bigr)
\]
be \emph{any} linear map into the $\gamma_n$-almost-everywhere equivalence
classes of measurable $\mathbb{R}^n$-valued functions such that
\begin{enumerate}
\item[(T1)] $T(\iota v) = R_\gamma v$ for every $v \in L^2(\gamma_n)$, with
$\iota$ the inclusion of \noderef{ToL1} and $R_\gamma$ the transform of
\noderef{RieszTransform}; and
\item[(T2)] $t\,\gamma_n\bigl(\{x : |Th(x)|_{\mathbb{R}^n} > t\}\bigr)
\leq 2\,\|h\|_{L^1(\gamma_n)}$ for every $h \in L^1(\gamma_n;\mathbb{R})$ and
every real $t > 0$.
\end{enumerate}
For a $\gamma_n$-integrable $h$, a coordinate $j$, a point $x$ and a real
$\varepsilon > 0$ write
$T_{j,\varepsilon}h(x) = \int_{|x-y|>\varepsilon}K_{\gamma,j}(x,y)h(y)\,dy$,
the $y$-integral against $\lambda$; for an element of $L^1(\gamma_n)$ this is
read through its canonical representative.

Let $f \in L^1(\gamma_n)$ and let $j \in \{1,\ldots,n\}$.  Then for
$\gamma_n$-almost every $x \in \mathbb{R}^n$ there is a real number $L$ with
\[
  \lim_{\varepsilon\downarrow0}T_{j,\varepsilon}f(x) \;=\; L ,
\]
the limit taken through $\varepsilon \in (0,\infty)$ as $\varepsilon \to 0$.

\emph{Existence only, and deliberately so.}  The statement asserts that the
limit exists; it says nothing about its value, and in particular does not
identify it with any coordinate of $Tf$.  That identification is a separate
assertion requiring separate inputs, and keeping it out of this statement is
what makes the two independently usable: a consumer that wants the value
obtains $L$ here and identifies it elsewhere.  Correspondingly no function
$x \mapsto L(x)$ is produced: the conclusion is an existential at each good
point, so no choice is made and no total function enters the interface.

\emph{Why the conclusion is almost-everywhere.}  The exceptional set is
genuine and not an artefact.  It is the set where the oscillation of
$\varepsilon \mapsto T_{j,\varepsilon}f(x)$ at $0$ fails to vanish, and it is
controlled only in measure, through the weak-type bound of
\noderef{GaussMaximalTruncationLocal} applied to the errors of a sequence of
smooth approximants.  By contrast each individual $T_{j,\varepsilon}f(x)$ is a
finite real number at \emph{every} $x$, by \noderef{GaussRieszTruncationSmoothSplit};
it is only the limit that needs an exceptional set.

\emph{Both (T1) and (T2) are assumed, because the input carries both.}  The
weak-type bound of \noderef{GaussMaximalTruncationLocal} is a statement about
every $T$ satisfying the two, so this lemma is too.  It is not vacuous because
\noderef{L1Extension} exhibits such a $T$.
\end{lemma}
\begin{proof}
\emph{Step 1: the approximants, their errors, and the splitting.}
For each $k \in \mathbb{N}$ apply \noderef{GaussRieszTruncationSmoothSplit} to
$f$ and $k$.  It produces a function $\varphi_k$ and an element
$g_k \in L^1(\gamma_n)$ such that $\varphi_k$ is smooth with compact support
and $\gamma_n$-integrable, such that
\begin{equation*}
  \int_{\mathbb{R}^n}|g_k|\,d\gamma_n < \frac1{k+1} , \tag{1}
\end{equation*}
and such that for every $j$, every $x$ and every real $\varepsilon > 0$ the
three integrands are Lebesgue integrable on $\{y : |x-y| > \varepsilon\}$ and
\begin{equation*}
  T_{j,\varepsilon}f(x)
    = T_{j,\varepsilon}\varphi_k(x) + T_{j,\varepsilon}g_k(x) . \tag{2}
\end{equation*}
All quantities appearing in $(2)$ are finite real numbers, at every $x$ and
every $\varepsilon > 0$.

\emph{Step 2: off a $\gamma_n$-null set, the errors' maximal truncations can be
made arbitrarily small.}  We claim that for $\gamma_n$-almost every
$x \in \mathbb{R}^n$,
\begin{equation*}
  \text{for every } m \in \mathbb{N} \text{ there is } k \in \mathbb{N}
  \text{ with } \bigl|T_{j,\varepsilon}g_k(x)\bigr| \leq \frac1{m+1}
  \text{ for every real } \varepsilon > 0 . \tag{3}
\end{equation*}

Let $B$ be the set of $x$ at which $(3)$ fails.  If $x \in B$ there is
$m \in \mathbb{N}$ such that for every $k$ the bound fails, i.e.\ for every $k$
there is a real $\varepsilon > 0$ with
$|T_{j,\varepsilon}g_k(x)| > 1/(m+1)$; and since
$\mathrm{ofReal}$ is strictly increasing on $[0,\infty)$ and a single term is
at most the supremum,
\[
  \mathrm{ofReal}\Bigl(\frac1{m+1}\Bigr)
    \;<\; \sup_{\varepsilon>0}\mathrm{ofReal}\bigl|T_{j,\varepsilon}g_k(x)\bigr|
  \qquad\text{for every } k ,
\]
the supremum taken in $[0,+\infty]$.  Also $|x| \leq r+1$ for some
$r \in \mathbb{N}$, for instance $r = \lceil|x|\rceil$.  Hence
\[
  B \;\subseteq\; \bigcup_{m\in\mathbb{N}}\bigcup_{r\in\mathbb{N}}A_{r,m},
  \qquad
  A_{r,m} = \Bigl\{x : |x| \leq r+1 \text{ and }
    \mathrm{ofReal}\bigl(\tfrac1{m+1}\bigr)
      < \sup_{\varepsilon>0}\mathrm{ofReal}\bigl|T_{j,\varepsilon}g_i(x)\bigr|
    \text{ for every } i\Bigr\} .
\]
By countable subadditivity of the outer measure $\gamma_n$ it therefore
suffices to show $\gamma_n(A_{r,m}) = 0$ for each fixed $r$ and $m$.

Fix $r$ and $m$.  Apply \noderef{GaussMaximalTruncationLocal} to this $T$ with
its two properties, which are (T1) and (T2) and hold by assumption, to the
coordinate $j$, and to the radius $R := r+1$, whose hypothesis $R \geq 1$ holds
because $r \geq 0$.  It supplies a constant
$C_r \in (0,\infty)$ --- in particular $C_r \neq 0$ and $C_r \neq +\infty$, and
produced before its input is quantified, hence the same for every input ---
such that for every $h \in L^1(\gamma_n)$ and every $t \in (0,+\infty]$,
\begin{equation*}
  \gamma_n\Bigl(\Bigl\{x : |x| \leq r+1 \text{ and }
    t < \sup_{\varepsilon>0}\mathrm{ofReal}\bigl|T_{j,\varepsilon}h(x)\bigr|
  \Bigr\}\Bigr)
  \;\leq\; \frac{C_r\cdot\mathrm{ofReal}\bigl(\int|h|\,d\gamma_n\bigr)}{t} .
  \tag{4}
\end{equation*}
Put $t_m = \mathrm{ofReal}(1/(m+1))$, which is nonzero because
$1/(m+1) > 0$, and $D = C_r/t_m$, which is not $+\infty$ because $C_r$ is not
and $t_m$ is not $0$.  For each $k$ we have
$A_{r,m} \subseteq \{x : |x| \leq r+1 \text{ and } t_m <
 \sup_{\varepsilon>0}\mathrm{ofReal}|T_{j,\varepsilon}g_k(x)|\}$, by taking
$i = k$ in the defining condition.  So by monotonicity, then $(4)$ with
$h = g_k$ and $t = t_m$, then $(1)$ and the monotonicity of $\mathrm{ofReal}$
and of multiplication in $[0,+\infty]$,
\[
  \gamma_n(A_{r,m})
    \;\leq\; \frac{C_r\cdot\mathrm{ofReal}\bigl(\int|g_k|\,d\gamma_n\bigr)}{t_m}
    \;=\; D\cdot\mathrm{ofReal}\Bigl(\int|g_k|\,d\gamma_n\Bigr)
    \;\leq\; D\cdot\mathrm{ofReal}\Bigl(\frac1{k+1}\Bigr) ,
\]
the middle equality being $ab/c = (a/c)b$ in $[0,+\infty]$.  This holds for
every $k$.  Now $1/(k+1) \to 0$ as $k \to \infty$, so
$\mathrm{ofReal}(1/(k+1)) \to 0$ in $[0,+\infty]$, and since $D \neq +\infty$
also $D\cdot\mathrm{ofReal}(1/(k+1)) \to D\cdot 0 = 0$.  A quantity in
$[0,+\infty]$ bounded above by every term of a sequence converging to $0$ is
at most $0$, so $\gamma_n(A_{r,m}) = 0$, proving $(3)$.

\emph{Step 3: at a point where $(3)$ holds the truncations are Cauchy.}
Let $x$ be such a point.  We claim that for every real $\eta > 0$ there is a
real $\rho > 0$ with
\begin{equation*}
  \bigl|T_{j,\varepsilon}f(x) - T_{j,\delta}f(x)\bigr| \leq \eta
  \qquad\text{for all }\varepsilon,\delta \in (0,\rho) . \tag{5}
\end{equation*}

Fix $\eta > 0$.  Choose $m \in \mathbb{N}$ with $1/(m+1) < \eta/8$, which is
possible since $\eta/8 > 0$, and let $k$ be as in $(3)$ for this $m$, so
\begin{equation*}
  \bigl|T_{j,\varepsilon}g_k(x)\bigr| \leq \frac1{m+1}
  \qquad\text{for every real }\varepsilon > 0 . \tag{6}
\end{equation*}
Since $\varphi_k$ is smooth with compact support, \noderef{PVCoreSmooth}
applies to it at the coordinate $j$ and the point $x$, and its clause (2) gives
a real number $\Phi$ with
$T_{j,\varepsilon}\varphi_k(x) \to \Phi$ as $\varepsilon \downarrow 0$.  By the
definition of that limit there is a real $\rho > 0$ such that
\begin{equation*}
  \bigl|T_{j,\varepsilon}\varphi_k(x) - \Phi\bigr| < \frac\eta8
  \qquad\text{for every }\varepsilon \in (0,\rho) . \tag{7}
\end{equation*}
Let $\varepsilon,\delta \in (0,\rho)$.  Writing
$a = T_{j,\varepsilon}\varphi_k(x)$, $b = T_{j,\varepsilon}g_k(x)$,
$c = T_{j,\delta}\varphi_k(x)$, $d = T_{j,\delta}g_k(x)$, all finite reals by
Step 1, the identity
$a + b - (c+d) = \bigl((a-\Phi) + (\Phi-c)\bigr) + \bigl(b + (-d)\bigr)$ and
the triangle inequality for the absolute value give
\[
  |a + b - (c+d)| \;\leq\; |a-\Phi| + |c-\Phi| + |b| + |d| ,
\]
using $|\Phi - c| = |c - \Phi|$ and $|-d| = |d|$.  By $(2)$ at $\varepsilon$
and at $\delta$ the left-hand side is
$|T_{j,\varepsilon}f(x) - T_{j,\delta}f(x)|$, and by $(7)$ and $(6)$ the
right-hand side is at most
\[
  \frac\eta8 + \frac\eta8 + \frac1{m+1} + \frac1{m+1}
  \;<\; \frac\eta8 + \frac\eta8 + \frac\eta8 + \frac\eta8
  \;=\; \frac\eta2 \;\leq\; \eta .
\]
This is $(5)$.

\emph{Step 4: a Cauchy family in a complete space converges.}
Still with $x$ as in Step 3, let $\mathcal{F}$ be the filter of
neighbourhoods of $0$ relative to $(0,\infty)$, i.e.\ the filter generated by
the sets $(0,\rho)$ for $\rho > 0$; it is not the trivial filter, because every
such set is nonempty.  Let $\mathcal{G}$ be the image of $\mathcal{F}$ under
$\varepsilon \mapsto T_{j,\varepsilon}f(x)$, a filter on $\mathbb{R}$, which is
again nontrivial.  We check that $\mathcal{G}$ is Cauchy, i.e.\ that for every
real $\eta > 0$ there is a set $S \in \mathcal{G}$ all of whose pairs of points
are at distance less than $\eta$.  Given $\eta > 0$, apply $(5)$ with $\eta/2$
to get $\rho > 0$, and take
$S = \{T_{j,\varepsilon}f(x) : \varepsilon \in (0,\rho)\}$.  Then
$S \in \mathcal{G}$, because $(0,\rho)$ belongs to $\mathcal{F}$ and is
contained in the preimage of $S$; and any two points of $S$ are
$T_{j,\varepsilon}f(x)$ and $T_{j,\delta}f(x)$ for some
$\varepsilon,\delta \in (0,\rho)$, hence at distance at most $\eta/2 < \eta$ by
$(5)$.  Since $\mathbb{R}$ is complete (Preamble), the Cauchy filter
$\mathcal{G}$ converges: there is a real $L$ with $\mathcal{G}$ finer than the
neighbourhood filter of $L$, which is exactly the assertion that
$T_{j,\varepsilon}f(x) \to L$ as $\varepsilon \downarrow 0$.

As the set of $x$ for which $(3)$ holds is of full $\gamma_n$-measure by Step 2,
this completes the proof.
\end{proof}
